\documentclass{beamer}
\usetheme{Berkeley}

\usepackage{
  amsmath,mathtools,amssymb,amsthm,
  centernot,float,diagbox,adjustbox,
  graphicx,caption,subcaption,listings,
  enumerate,xcolor,hyperref,color,soul,
  stmaryrd,amsrefs,tikz,mathrsfs,etoolbox,
  refcount,multicol,MnSymbol,yfonts,framed,
  relsize,colortbl,pgf-umlsd,varwidth,
  environ,xparse,forest}
\usepackage[many]{tcolorbox}
\forestset{qtree/.style={for tree={parent anchor=south, 
           child anchor=north,align=center,inner sep=0pt}}}
\usepackage[framemethod=tikz]{mdframed}
\newcommand{\grayleftbar}[1]{\textcolor[gray]{0.7}{%
     \begin{leftbar}\textcolor{black}{#1}\end{leftbar}}}   
\newcommand{\darkleftbar}[1]{\textcolor[gray]{0.3}{%
     \begin{leftbar}\textcolor{black}{#1}\end{leftbar}}}   
\newcounter{countitems}
\newcounter{nextitemizecount}
\newcommand{\setupcountitems}{%
  \stepcounter{nextitemizecount}%
  \setcounter{countitems}{0}%
  \preto\item{\stepcounter{countitems}}%
}
\makeatletter
\newcommand{\computecountitems}{%
  \edef\@currentlabel{\number\c@countitems}%
  \label{countitems@\number\numexpr\value{nextitemizecount}-1\relax}%
}
\newcommand{\nextitemizecount}{%
  \getrefnumber{countitems@\number\c@nextitemizecount}%
}
\newcommand{\previtemizecount}{%
  \getrefnumber{countitems@\number\numexpr\value{nextitemizecount}-1\relax}%
}

\usepackage[utf8]{inputenc}
\makeatother    
\newenvironment{AutoMultiColItemize}{%
\ifnumcomp{\nextitemizecount}{>}{3}{\begin{multicols}{2}}{}%
\setupcountitems\begin{itemize}}%
{\end{itemize}%
\unskip\computecountitems\ifnumcomp{\previtemizecount}{>}{3}{\end{multicols}}{}}
\usetikzlibrary{
    arrows,shapes,positioning,decorations.markings,
    decorations.pathreplacing,hobby,automata,tikzmark}
\definecolor{mycolor}{rgb}{0.122, 0.435, 0.698}

\newmdenv[innerlinewidth=0.5pt, roundcorner=4pt,linecolor=mycolor,innerleftmargin=6pt,
innerrightmargin=6pt,innertopmargin=6pt,innerbottommargin=6pt]{mybox}

\usepackage[pdf]{graphviz}
\usetikzlibrary{shapes,arrows,calc}
\usepackage[utf8]{inputenc}
\newtheorem{remark}{Remark}

\newcommand{\TODO}[0]{\textcolor{red}{TODO}}
\newcommand{\F}[0]{\mathsf{F}}
\newcommand{\U}[0]{\mathsf{U}}
\newcommand{\G}[0]{\mathsf{G}}
\newcommand{\X}[0]{\mathsf{X}}
\newcommand{\acs}[0]{\textsc{ACL2S}}
\newcommand{\acl}[0]{\textsc{ACL2}}
\newcommand{\pr}[0]{\textsc{Promela}}
\newcommand{\ko}[0]{\textsc{Korg}}
\newcommand{\spn}[0]{\textsc{SPIN}}
\newcommand{\ie}{\emph{i.e.}}
\newcommand{\etal}{\emph{et al.}}

\newcommand{\Pcl}[0]{\textsc{P}}
\newcommand{\NP}[0]{\textsc{NP}}
\newcommand{\SAT}[0]{\texttt{SAT}}
\newcommand{\coNP}[0]{\textsc{coNP}}
\newcommand{\Time}[0]{\text{TIME}}
\newcommand{\Space}[0]{\text{SPACE}}
\newcommand{\threeSAT}[0]{\texttt{3SAT}}

\DeclarePairedDelimiter\abs{\lvert}{\rvert}
\DeclarePairedDelimiter\norm{\lVert}{\rVert}

\usepackage{tikz}
\usetikzlibrary{arrows,chains,matrix,positioning,scopes}

\newlength{\bubblewidth}
\AtBeginDocument{\setlength{\bubblewidth}{.55\textwidth}}
\definecolor{bubblegreen}{RGB}{103,184,104}
\definecolor{bubblegray}{RGB}{241,240,240}
\newcommand{\bubble}[4]{%
  \tcbox[
    colback=#1,
    colframe=#1,
    #2,
  ]{\color{#3}\begin{varwidth}{\bubblewidth}#4\end{varwidth}}%
}

\ExplSyntaxOn
\seq_new:N \l__ooker_bubbles_seq
\tl_new:N \l__ooker_bubbles_first_tl
\tl_new:N \l__ooker_bubbles_last_tl

\NewEnviron{rightbubbles}
 {
  \raggedleft\sffamily
  \seq_set_split:NnV \l__ooker_bubbles_seq { \par } \BODY
  \int_compare:nTF { \seq_count:N \l__ooker_bubbles_seq < 2 }
   {
    \bubble{bubblegreen}{rounded~corners}{white}{\BODY}
   }
   {
    \seq_pop_left:NN \l__ooker_bubbles_seq \l__ooker_bubbles_first_tl
    \seq_pop_right:NN \l__ooker_bubbles_seq \l__ooker_bubbles_last_tl
    \bubble{bubblegreen}{sharp~corners=southeast}{white}{\l__ooker_bubbles_first_tl}\par
    \seq_map_inline:Nn \l__ooker_bubbles_seq
     {
      \bubble{bubblegreen}{sharp~corners=east}{white}{##1}\par
     }
    \bubble{bubblegreen}{sharp~corners=northeast}{white}{\l__ooker_bubbles_last_tl}\par
   }
 }
\NewEnviron{leftbubbles}
 {
  \raggedright\sffamily
  \seq_set_split:NnV \l__ooker_bubbles_seq { \par } \BODY
  \int_compare:nTF { \seq_count:N \l__ooker_bubbles_seq < 2 }
   {
    \bubble{bubblegray}{rounded~corners}{black}{\BODY}
   }
   {
    \seq_pop_left:NN \l__ooker_bubbles_seq \l__ooker_bubbles_first_tl
    \seq_pop_right:NN \l__ooker_bubbles_seq \l__ooker_bubbles_last_tl
    \bubble{bubblegray}{sharp~corners=southwest}{black}{\l__ooker_bubbles_first_tl}\par
    \seq_map_inline:Nn \l__ooker_bubbles_seq
     {
      \bubble{bubblegray}{sharp~corners=west}{black}{##1}\par
     }
    \bubble{bubblegray}{sharp~corners=northwest}{black}{\l__ooker_bubbles_last_tl}\par
   }
 }
\ExplSyntaxOff


\title{Communication Complexity}
\author{Jake Kinsella and Max von Hippel}
\institute{Northeastern University}
\date{\today}



\begin{document}

\frame{\titlepage}

\begin{frame}
\frametitle{Communication Complexity}
\emph{If Alice knows $x$, and Bob knows $y$, how many bits of information must they communicate, in order for both Alice and Bob to know $f(x, y)$?}
\end{frame}

\begin{frame}
\tableofcontents
\end{frame}

% 0  - 5m   - Max introduces the problem
% 5  - 10m  - Jake a couple examples
% 10 - 30m  - Jake does fooling set
% 30 - 50m  - Max does the tiling method
% 50 - 60m  - Max does 2-party discrepancy method
% 60 - 65m  - Jake introduces motivating problem for multi-party version
% 65 - 70m  - Jake defines generalizes multi-party version of problem 
% 70 - 85m  - Max gives multi-party discrepancy without proof
% 85 - 90m  - Jake does non-deterministic version
% 90 - 95m  - Max does randomized

\section{Introduction}

\begin{frame}{(Max)}
Consider a two-party communication problem,
	in which the participants

\begin{figure}[h]
\centering
\begin{subfigure}{.3\textwidth}
  \centering
  \includegraphics[width=.2\linewidth]{introduction/alice.png}
  \caption{Alice}
  \label{fig:alice}
\end{subfigure}%
and
\begin{subfigure}{.3\textwidth}
  \centering
  \includegraphics[width=.2\linewidth]{introduction/bob.png}
  \caption{Bob}
  \label{fig:bob}
\end{subfigure}
\label{fig:participants}
\end{figure}

\emph{participate} to compute a function:
\[f : \underbrace{\mathbb{B}^n}_{\begin{aligned}\text{Alice's}\\\text{input}\end{aligned}} \times 
      \underbrace{\mathbb{B}^n}_{\begin{aligned}\text{Bob's}\\\text{input}\end{aligned}} \to 
      \underbrace{\mathbb{B}}_{\begin{aligned}\text{global}\\\text{output}\end{aligned}}\]
\end{frame}

\begin{frame}{(Max)}
The players can come up with a \emph{protocol} 
\(\Pi = (p_1, ..., p_t)\),
namely, 
	for some natural $t \in \mathbb{N}$,
	a sequence of $t$-many functions
\(p_i : \mathbb{B}^* \to \mathbb{B}^*\)
such that the communication between the players looks like this ...
\end{frame}

\begin{frame}{(Max)}
{\color{blue}{\footnotesize Alice is given input $x$.}}
\pause
\begin{leftbubbles}
{\footnotesize Hello Bob.  I can't reveal $x$, but $p_1(x)$ is \texttt{p1}.}
\end{leftbubbles}
\pause
\rightline{{\color{olive}{\footnotesize Bob is given input $y$.}}}
\pause
\begin{rightbubbles}
{\footnotesize Thanks Alice.  I can't reveal $y$, but $p_2(y, \texttt{p1})$ is \texttt{p2}.}
\end{rightbubbles}
\pause
\center{\footnotesize{... yada yada yada ...}}
\pause
\begin{leftbubbles}
{\footnotesize Pleasure doing business with you Bob.  My final clue for you is that $p_{n-1}(x, \texttt{p1}, ..., \texttt{pn-2})$ is \texttt{pn-1}.}
\end{leftbubbles}
\pause
\begin{rightbubbles}
{\footnotesize Rad.  Then $p_n(y, \texttt{p1}, ..., \texttt{pn-1})$ is \texttt{pn}.  TTFN!}
\end{rightbubbles}
\end{frame}

\begin{frame}{(Max)}
\begin{itemize}
\item The functions $p_i$ can be \emph{anything} so long as they are well-defined.  E.g., could solve the Halting Problem.
\item After the final message, \emph{both parties} must know $f(x, y)$.
\end{itemize}
\pause
\begin{definition}[Communication Complexity]
Suppose $\Pi$ is a protocol for $f$ in which at most $t$
bits are communicated between Alice and Bob.
Then the \emph{communication complexity} of $\Pi$,
	denoted $C(\Pi)$, is $t$.
\end{definition}
\pause
\begin{definition}[$C(f)$]
The communication complexity of $f$, denoted $C(f)$,
	is the minimum communication complexity
	achieved by any protocol for $f$.
\end{definition}
\end{frame}

\subsection{Examples}

\begin{frame}{Parity (Jake)}
\begin{example}[Are the number of \texttt{1}s in $xy$ even (0), or odd (1)?]
\(f : \mathbb{B}^n \times \mathbb{B}^n \to \mathbb{B}\)
is precisely
\((x, y) \mapsto \bigoplus xy\).
\end{example}
\pause
Example protocol $\Pi$:
\begin{leftbubbles}
{\footnotesize $\texttt{P1}=\text{parity}(x)$.}
\end{leftbubbles}
\pause
\begin{rightbubbles}
{\footnotesize $\texttt{P2}=\text{parity}(y) \bigoplus \texttt{P1}$}
\end{rightbubbles}
\pause
Now both Alice and Bob know $f(x, y) = \texttt{P2}$.
$C(f) \leq 2$ because $C(\Pi) = 2$ and $\Pi$ implements $f$.
But $C(f) \geq 2$ because $f$ depends on $x$ and $y$.
Hence \(C(f) = 2\).
\end{frame}

\begin{frame}{Halting (Jake)}
\begin{example}[$A_{TM}$]
\(H : \mathbb{B}^n \times \mathbb{B}^n \to \mathbb{B}\)
is precisely
\(\langle M, 1^{n} \rangle \mapsto 1\)
if $M$ halts on $1^{n}$ else 0.
\end{example}
\pause
Example protocol $\Pi$:\\
\begin{rightbubbles}
{\footnotesize $\texttt{P1}=1 \text{ if } y = 1^{n} \text{ else } 0$.}
\end{rightbubbles}
\pause
\begin{leftbubbles}
\texttt{P2}$=(M$ does/doesn't halt on $1^{n})$.
\end{leftbubbles}
\pause
Both players have unlimited computation power.
We are only interested in communication complexity.
\end{frame}

\section{Methods}

\subsection{2-Party Problem}

\begin{frame}{Methods (Max)}
If we find a protocol $\Pi$,
	then we know $C(f)$ is at most $C(\Pi)$.
\pause

What if we don't know any protocol $\Pi$?
\pause
\begin{itemize}
	\item Could we upper-bound $C(f)$ without knowing $\Pi$?
\end{itemize}

\pause
What if the only protocols we find seem really lousy?
\pause
\begin{itemize}
	\item Could we lower-bound $C(f)$ without finding a better protocol?
\end{itemize}

\pause
\centering{\textbf{TL;DR: yup.}}
\end{frame}

\subsubsection{Fooling Set Method}

\begin{frame}{Fooling Set Method (Jake)}
We begin with a motivating observation.
\pause
\begin{lemma}[Communication Equality is Image Equality]
If Alice and Bob exchange the same sequence of messages when Alice gets $\textbf{x}$ and Bob gets $\textbf{y}$ as they do when Alice gets $\textbf{x}'$ and Bob gets $\textbf{y}'$, then $f(\textbf{x}, \textbf{y}) = f(\textbf{x}', \textbf{y}')$.
\end{lemma}
\pause
\begin{proof}
$\Pi$ is deterministic and $f$ is a function.
\end{proof}
\pause
\underline{Idea:} an efficient protocol will efficiently group together inputs that go to the same output.

\end{frame}

\begin{frame}{Fooling Set Method (Jake)}
\begin{definition}[Fooling Set]
If $f : \mathbb{B}^n \times \mathbb{B}^n \to \mathbb{B}$ is a function,
	a \emph{fooling set} for $f$ is a set
	\(S \subseteq \mathbb{B}^n \times \mathbb{B}^n\)
	such that for some choice $b \in \mathbb{B}$
\(f(S) = \{ b \}\)
but, for all distinct $(x, y), (x', y') \in S$,
either $f(x, y') \neq b \text{ or } f(x', y) \neq b$
\end{definition}
Basically, a fooling set is a group of inputs that go to the same output,
but which is \emph{brittle} to argument-swapping.
In some sense these \emph{brittle} sets lower-bound the difficulty in grouping like inputs.
\pause
\begin{lemma}[Fooling Set Method]
If $f$ has a size-$M$ fooling set, then $C(f) \geq \text{log}_2(M)$.
\end{lemma}
\end{frame}

\begin{frame}{Fooling Set Method (Jake)}
\begin{example}[Set-Disjointness]
$\textsc{Disj} : \mathbb{B}^n \times \mathbb{B}^n \to \mathbb{B}$ is the function
that maps $(A, B)$ to 1 if $A \cap B = \emptyset$ else 0.
\end{example}

How many fooling sets does \textsc{Disj} have?
\pause
Notice $A \cap B = \emptyset$ if $B = \overline{A}$.
\pause
There are $2^n$ possible values $A$.
\pause
None of these distinct $(A, \overline{A}), (A', \overline{A'})$ satisfy
 $A \cap \overline{A} = A \cap \overline{A'}$ or
$A \cap \overline{A} = A' \cap \overline{A}$
else they wouldn't be distinct.
\pause
So we get a $2^n$-size fooling set.
\pause
\[\therefore \, C(\textsc{Disj}) \geq \text{log}_2(2^n) = n\]
\end{frame}

\begin{frame}{Fooling Set Method (Jake)}
\textbf{NTS:} If $f$ has a size-$M$ fooling set then $C(f) \geq \text{log}_2(M)$.
\begin{proof}
For a contradiction suppose a protocol $\Pi$ exists for $f$ s.t. $C(\Pi) < \text{log}_2(M)$.
\pause
Then $\Pi$ yields at most $2^{C(\Pi)} < 2^{\text{log}_2(M)} = M$
distinct communication patterns.
\pause
However, there are $M$ input pairs $(x, y)$ that each have distinct communication patterns by definition.
\pause
Since $2^{C(\Pi)} < M$ there must be some $(x, y), (x', y')$ on which $\Pi$ yields the same communication pattern.

\medskip

\pause
Then $(x, y')$ must yield the same communication pattern as $(x, y)$
as Bob cannot possibly tell the difference.  The argument is symmetric for $(x', y)$ and Alice.  One of the two must yield a contradiction and we are done.
\end{proof}
\end{frame}

\subsubsection{Tiling Method}

\begin{frame}{Tiling Method (Max)}
\TODO
\end{frame}

\subsubsection{Discrepency Method}

\begin{frame}{2-Party Discrepency Method (Max)}
\TODO
\end{frame}

\subsection{Multi-Party Problem}

\begin{frame}{Multi-Party Problem (Jake)}
\TODO
\end{frame}

\subsubsection{Multi-Party Discrepency Method}

\begin{frame}{Multi-Party Discrepency Method (Max)}
\TODO
\end{frame}

\section{Other Variants}

\subsection{Non-Deterministic}

\begin{frame}{Non-Deterministic (Jake)}
\TODO
\end{frame}

\subsection{Randomized}

\begin{frame}{Randomized (Max)}
\TODO
\end{frame}

\end{document}